\scanpage{2}
\margin{Pise, Mai 1966}
\mspage{1}

Cher John,

J'ai réfléchi aux groupes formels et à la cohomologie de De Rahm, et suis arrivé à un projet de théorie, ou plutôt \inserted{de} début de théorie, que j'ai envie de t'exposer,pour me clarifier les idées.

\uline{Chapitre} 1\qquad\uline{La notion de cristal.}

Commentaire terminologique: Un cristal possède deux propriétés caractéristiques: la \uline{rigidité}, et la faculté de \uline{croitre},dans un voisinage approprié. Il y a des cristaux de toute espèce de substance: des cristaux de soude, de souf\deleted{f}re, de modules, d'anneaux, de schémas relatifs etc.

1.0. Soit S un préschéma, au dessus d'un autre R; dans le cas qui nous intéressera le plus, on aura $R=\operatorname{Spec}(\underline Z)$. Soit C la catégorie des \inserted{R-}préschémas \inserted{T} \uline{sous} S, \inserted{(}i.e. munis d'un R-\deleted{\unsure{ho}}morphisme $S\to T$,\inserted{)} tels que $S\to T$ soit une immersion fermée définie par un idéal localement nilpotent. On a le foncteur ``oubli'' de C dans (Sch), et la catégorie fibrée des Modules quasi-cohérents sur des préschémas variables induit donc une catégorie fibrée sur C, associant à tout T la catégorie des Modules quasi-cohérents sur T.

\uline{Définition} 1.1. On appelle cristal de modules (sous-entendu: quasi-cohérents) sur S, relativement à R, une section cartésienne de la catégorie fibrée précédente au dessus de C. Plus généralement, pour toute catégorie fibrée $\underline F$ sur $(\mathrm{Sch})_{/R}$, on définit la notion de ``$\underline F$-cristal'' sur S, ou ``cristal en objets de $\underline F$ '', de la façon correspondante. Ceci donne un sens aux expressions: cristal en algèbres, en algèbres commutatives, en préschémas relatifs etc\inserted{,} sur S relativement à R. Quand $R=\operatorname{Spec}(\underline Z)$, on parlera de ``cristal absolu'' sur S, de l'espèce considérée.

1.2. Les $\underline F$-cristaux sur S forment une catégorie, qui dépend fonctoriellement de $\underline F$. Ceci permet, en particulier, de définir sur la catégorie des cristaux de modules sur S les opérations tensorielles habituelles: produit tensoriel de deux cristaux de modules etc, satisfaisant aux propriétés habituelles.
\margin{1}
